\begin{recipe}
[% 
    preparationtime = {\unit[8]{min}},
    portion = {\portion{4}},
    source = {\protect\recipesource{https://www.coffeecircle.com/de/e/french-press-anleitung}{Coffee Circle: French Press Anleitung}},
    calory = {2kcal | 0g Protein | 0g KH | 0g Fett},
    price = {0,40€},
    created = {10.06.2026},
    %rating = {1},
]
{French Press Kaffee}

%\graph{% pictures
%    big=pic/getraenke/03_french-press
%}

\introduction{%
French Press ist Kaffee ohne großes Theater: Kaffee rein, Wasser drauf, warten, runterdrücken. Durch das Metallsieb bleibt mehr Körper in der Tasse, also eher kräftig und voll statt super clean.
}

\ingredients{
    65 g & Kaffee, mittel-grob gemahlen\index{Getränke!Kaffee}\\
    1 l & Wasser
}

\preparation{%
    \step%
    Kaffee frisch mittel-grob mahlen.
    
    \step%
    French Press optional mit heißem Wasser vorwärmen und das Wasser wieder ausschütten.
    
    \step%
    Kaffeepulver in die French Press geben und mit 1 l heißem Wasser bei etwa 96 °C zügig aufgießen.
    
    \step%
    Kurz umrühren, damit das Kaffeepulver vollständig mit Wasser benetzt ist.
    
    \step%
    Deckel aufsetzen. Den Stempel nur so weit nach unten drücken, dass das Kaffeepulver unter Wasser bleibt.
    
    \step%
    4 Minuten ziehen lassen.
    
    \step%
    Nach etwa 3{,}5 Minuten zweimal umrühren, um den Schaum zu entfernen. Danach den Stempel langsam und gleichmäßig nach unten drücken.
    
    \step%
    Direkt einschenken und trinken.
}

\hint{
    Kaffee nach dem Runterdrücken nicht ewig in der French Press stehen lassen, sonst extrahiert er weiter und wird bitter. Also entweder direkt austrinken oder umfüllen.
}

\suggestion[Anpassen]{%
    Wenn der Kaffee zu bitter wird, gröber mahlen oder etwas kürzer ziehen lassen. Wenn er dünn und wässrig schmeckt, feiner mahlen oder etwas mehr Kaffee nehmen.
}

\end{recipe}
